\begin{claim}{C.33}{Proved}{All proper operator groups on an open orbit}
Let $O$ be a connected noncompact second-countable surface without
boundary. Let $Q$ be a closed locally compact subgroup of
$\operatorname{Homeo}(O)$, with its compact-open topology, acting
freely and properly. Then inclusion induces an injection
\[
 \pi_1\!\left(C_{\operatorname{Homeo}(O)}(Q),\id\right)
       \longrightarrow \pi_1(\operatorname{Homeo}(O),\id).
                                                        \tag{C.33.1}
\]
For every open orbit of the canonically lifted hyperbolic
centralizer loop in C.28, its restriction therefore has a based
null-homotopy commuting with the entire restricted operator
stabilizer. This conclusion retains the restriction kernel in
C.30 and imposes no finiteness or discreteness hypothesis on the
original operator group.
It supplies an intrinsic null-homotopy on each open orbit;
continuity across orbit frontiers remains unproved.
\end{claim}
\begin{proof}
Write $N=C_{\operatorname{Homeo}(O)}(Q)$.
The topological group $Q$ acts faithfully on the connected
three-manifold $O\times\R$ by
\[
                         q(x,s)=(qx,s).
                                                        \tag{C.33.2}
\]
This action is continuous: on any compact subset of $O\times\R$
the first coordinate is controlled on its compact projection to
$O$. Pardon proves that a locally compact group acting faithfully
on a connected three-manifold is a Lie group
\cite[Theorem~1.5 and the reduction preceding it]{C_Pardon2013}.
Thus $Q$ is a Lie group. Since it is a subspace of the
second-countable homeomorphism group, it is second countable.
Set $A=Q^0$ and $d=\dim A$.
Properness and freeness make $q\mapsto qx$ a proper embedding of
$A$ in $O$ for each $x$. Invariance of dimension gives $d\leq2$.
The case $d=0$ is precisely C.32.

If $d=2$, invariance of domain makes every $A$-orbit open.
Connectedness gives a single orbit, and evaluation $A\to O$
is a homeomorphism. For $q\in Q$, choose $a\in A$ with $qx=ax$.
Then $a^{-1}q$ fixes $x$, so freeness gives $q=a$ and $Q=A$.
Under $O\cong A$ with its left translation action, the elements of
$N$ are exactly the right translations. The correspondence
is an anti-isomorphism of topological groups, and evaluation
$N\to O$ is a homeomorphism. A loop in $N$ null-homotopic
in $\operatorname{Homeo}(O)$ has a null evaluation loop in $O$,
and hence is null-homotopic in $N$. This proves (C.33.1) in
dimension two.

It remains to treat $d=1$. The connected one-dimensional Lie
group $A$ is $\R$ or $\R/\Z$, written additively.
Its restricted action is still free and proper.
Palais' slice theorem gives a locally trivial principal bundle
\[
                 p_0:O\longrightarrow B=O/A.
                                                        \tag{C.33.3}
\]
Here all hypotheses of the theorem hold: the surface is completely
regular, $A$ is Lie, and a relatively compact neighborhood has
relatively compact transporter by properness, which is the
Cartan condition. A slice at a point has trivial stabilizer.
The resulting local trivialization is the map
$(a,s)\mapsto as$ from $A\times S$ onto the open saturated
neighborhood of that slice. This is the free-action specialization
of \cite[Theorem~2.3.3]{C_Palais1961}; its maps and inverse are
also given in \cite[Theorem~3.7]{C_KramerMuratGarcia2026}.

We verify the dimension and local topology of the base rather
than assuming differentiability of this action. The elementary
fact needed is that a metrizable space $S$ for which
$S\times(-1,1)$ is a surface is a one-manifold without boundary.
Indeed, projecting path-connected open neighborhoods in the
product shows that $S$ is locally path connected. It has no
isolated point, since an isolated point would give an open
interval in a surface. Every embedded open arc in $S$ is open:
its product with $(-1,1)$ maps injectively into a surface, and
invariance of domain applies.

Fix $p\in S$. Local path connectedness and the absence of isolated
points give an embedded arc $\alpha:[0,1]\to S$ with
$\alpha(0)=p$. Choose a path-connected open neighborhood $W$
of $p$ avoiding $\alpha(1)$. If $W$ contains a point outside
$\alpha([0,1])$, join $p$ to that point by an embedded arc in
$W$. This second arc cannot enter and subsequently leave
$\alpha((0,1))$: that open arc has boundary contained in
$\{p,\alpha(1)\}$, and an embedded arc in $W$ cannot revisit
$p$ or pass through $\alpha(1)$. Thus the two arcs meet only
at $p$, and together contain an arc having $p$ as an interior
point. Such an open arc is an open neighborhood of $p$ by the
preceding paragraph.
In the other case $W\subseteq\alpha([0,1))$, and a smaller
neighborhood of $p$ is a half interval. Its product with
$(-1,1)$ has zero local second homology at $(p,0)$,
whereas a surface without boundary has local second homology
$\Z$. This is impossible. The elementary fact follows.
Applied to the slices in (C.33.3), it proves that $B$ is a
one-manifold without boundary. Properness gives a Hausdorff
quotient, the open quotient map gives second countability, and
connectedness is inherited from $O$. Consequently $B$ is
$\R$ or $\R/\Z$.

Every principal $A$-bundle over either of these one-manifolds
is trivial. One can see the actual section as follows.
Over an interval, local trivializations are joined successively
across overlapping subintervals, adjusting one section by its
continuous $A$-valued transition function. A locally finite
interval exhaustion gives a section on $\R$.
Over the circle, first cut at one point and choose a section
on the resulting closed interval. Its endpoint discrepancy is
an element of the connected group $A$; multiplying the section
by a path from $0$ to the negative of that element makes the
endpoints agree, so that it descends to the circle.
The map $(a,b)\mapsto a\,s(b)$, with inverse given in each
bundle chart, therefore identifies
\[
                      O\cong A\times B,\qquad
             c(a,b)=(c+a,b)\quad(c\in A).                \tag{C.33.4}
\]
The pair $A=B=\R/\Z$ is excluded because $O$ is noncompact.

Let $\Delta=Q/A$. This discrete group acts freely and properly
on $B$. Freeness follows directly: an element of $Q$ fixing a
point of $B$ sends any point of its fibre to an $A$-translate;
after that translate is removed, it fixes a point of $O$ and
is the identity.
For properness, suppose $b_j\to b$ and
$\delta_j b_j\to b'$ in $B$. Local sections give $x_j\to x$
above $b_j$ and $y_j\to y$ above $\delta_j b_j$.
Choose a representative $q_j$ of $\delta_j$ and $a_j\in A$
such that $a_jq_jx_j=y_j$. Properness of the $Q$-action gives
a convergent subsequence of $a_jq_j$, so their components
$\delta_j$ are eventually constant on that subsequence.
This proves properness of the discrete action on the locally
compact metrizable space $B$.

We spell out this discrete action. On $\R$, an orientation
reverser has a fixed point and is excluded by freeness.
If $\Delta$ is nontrivial, its orbit of a point is closed and
discrete, with a nearest point to the right. The element
taking the chosen point to that next point generates
$\Delta$: its iterates are unbounded in both directions, and
any intervening orbit point would contradict the choice of
the nearest one. Choosing coordinates on one fundamental
interval and extending by its iterates identifies the action
with
\[
                   \Delta=\Z,\qquad \beta(b)=b+1.
                                                        \tag{C.33.5}
\]
On $\R/\Z$, properness makes $\Delta$ finite and freeness
again excludes orientation reversers. The action on one
finite orbit preserves its cyclic order, so it identifies
$\Delta$ with a cyclic group $C_n$. Coordinates on the
successive complementary intervals identify a generator
with
\[
                   \beta(b)=b+1/n\quad\text{in }\R/\Z.
                                                        \tag{C.33.6}
\]
The trivial group is also allowed on either base.

Suppose $\Delta$ is nontrivial and choose a representative
$q\in Q$ of the indicated generator. Conjugation by $q$ on
$A$ is multiplication by $\lambda$, where
$\lambda\in\R^\times$ for $A=\R$, and
$\lambda\in\{1,-1\}$ for $A=\R/\Z$.
Equation (C.33.4) therefore expresses its action exactly as
\[
                 q(a,b)=(\lambda a+\tau(b),\beta(b))
                                                        \tag{C.33.7}
\]
for a continuous function $\tau:B\to A$.
We next give a coordinate change removing this translation
term, with all component relations retained.

For $B=\R$ and $\Delta=\Z$, choose a continuous function
$c:[0,1]\to A$ whose endpoints satisfy
$c(1)=\lambda c(0)+\tau(0)$, and extend it in both directions
by
\[
                      c(b+1)=\lambda c(b)+\tau(b).
                                                        \tag{C.33.8}
\]
The endpoint equation makes the extension continuous at every
integer; invertibility of multiplication by $\lambda$ gives
the extension towards negative integers. In the new fibre
coordinate $a'=a-c(b)$, (C.33.7) becomes
$q(a',b)=(\lambda a',b+1)$.

For $B=\R/\Z$ and $\Delta=C_n$, necessarily $A=\R$.
Write $q^n=T_{a_0}$, where $T_{a_0}$ is translation by $a_0$
on the fibre. Conjugation gives $\lambda^n=1$ and
$\lambda a_0=a_0$.
If $\lambda=1$, replacing $q$ by $T_{-a_0/n}q$ makes
$q^n=1$ and does not change the group generated with $A$.
If $\lambda=-1$, then $n$ is even and $a_0=0$, so $q^n=1$
already. Thus the negative value of $\lambda$ is retained.
For this finite-order representative define the real-valued
function
\[
 c(b)=\frac1n\sum_{k=0}^{n-1}
      \operatorname{pr}_A\!\left(q^k(0,\beta^{-k}b)\right).
                                                        \tag{C.33.9}
\]
The affine nature of (C.33.7), followed by cyclic
reindexing using $q^n=1$, gives
$c(\beta b)=\lambda c(b)+\tau(b)$.
The same coordinate change $a'=a-c(b)$ now gives
$q(a',b)=(\lambda a',\beta b)$.
We have proved, by explicit changes of coordinates, that
the nonidentity components can be represented as
\[
                         q(a,b)=(\lambda a,\beta b).
                                                       \tag{C.33.10}
\]
When $\Delta$ is trivial no such additional equation is needed.

Commuting with all translations in (C.33.4) forces a
homeomorphism of $A\times B$ to have the form
\[
       h(a,b)=(a+\phi(b),\psi(b)),\qquad
       \phi\in C(B,A),\quad\psi\in\operatorname{Homeo}(B).
                                                       \tag{C.33.11}
\]
The inverse is
$(a,b)\mapsto(a-\phi(\psi^{-1}b),\psi^{-1}b)$.
These formulas and evaluation at the zero section give
mutually inverse continuous maps for the compact-open
topologies. In the nontrivial-component case, commuting
with (C.33.10) is exactly the pair of equations
\[
             \psi\beta=\beta\psi,\qquad
             \phi(\beta b)=\lambda\phi(b).
                                                       \tag{C.33.12}
\]
Thus the map (C.33.11) identifies $N$, as a space, with
the product of the function space defined by the second
equation and the homeomorphism centralizer defined by the
first. In particular the two constraints can be contracted
separately without changing the group $Q$ being centralized.

Let $h_t$ be a based loop in $N$ null-homotopic in
$\operatorname{Homeo}(O)$, and let $(\phi_t,\psi_t)$ be its
coordinates. The path $\psi_t$ preserves the orientation of
the one-manifold $B$.
When $A=\R$, multiplication of $\phi_t$ by $1-s$ contracts
the function-coordinate loop to zero and preserves its
equation in (C.33.12), including $\lambda=-1$.
For $B=\R$, C.27 contracts $\psi_t$ in the required
centralizer, with the trivial operator group when
$\Delta=1$. For $B=\R/\Z$, the evaluation loop
$h_t(0,b)=(\phi_t(b),\psi_t(b))$ is null-homotopic in
$\R\times(\R/\Z)$; hence $\psi_t(b)$ has winding zero.
C.29 contracts $\psi_t$ in the centralizer of the finite
rotation group in (C.33.6).

When $A=\R/\Z$, necessarily $B=\R$.
The based family $\phi_t(b)$ has a unique continuous lift
$\varphi:[0,1]\times\R\to\R$ with $\varphi_0=0$.
Evaluation of $h_t$ in the annulus $A\times\R$ shows that
$\varphi_1(b)=0$: this endpoint is the winding number of
the first coordinate of the null evaluation loop.
If $\Delta=\Z$, equation (C.33.12) says that
\[
           \varphi_t(b+1)-\lambda\varphi_t(b)\in\Z.
                                                       \tag{C.33.13}
\]
Its left side is a continuous function of $(t,b)$, equal
to zero at $t=0$, and is therefore identically zero.
Multiplying $\varphi$ by $1-s$ and descending to $\R/\Z$
contracts $\phi_t$ while preserving exact equivariance.
The same scaling applies when $\Delta=1$, without (C.33.13).
C.27 contracts the remaining $\psi_t$ on the line.
All these contractions keep $t=0,1$ fixed. Equations
(C.33.11)--(C.33.12), including the inverse formula, show
that their product is a continuous based contraction in
$N$. This completes the proof of (C.33.1).

Finally, for an open loop-group orbit from C.28, C.30
supplies the actual closed locally compact free proper
operator image $Q$. The restricted canonical plane loop
commutes with $Q$. It is null-homotopic in
$\operatorname{Homeo}(O)$ by the second assertion of
C.32, so (C.33.1) contracts it in the exact centralizer.
Every original operator in the setwise stabilizer acts
through its restriction to $Q$, and the restriction
kernel acts identically on $O$. The stated conclusion
follows with that kernel retained.
\end{proof}

\paragraph{What the open-orbit result provides.}
Equations (C.33.11)--(C.33.13) treat every component of the
proper operator group, including a circle fibre with a
nontrivial $\Z$-action on its line base. There is no
replacement of the base-translation centralizer by the
homeomorphism group of its quotient circle.
As in C.32, the obtained contraction can be transported to
other open orbits in the same operator orbit; exact
stabilizer equivariance makes the transport independent
of choices. This proof supplies no uniform control near
the boundary of the open orbit. C.34 subsequently proves
joint continuity, including inverses, on the entire closed
union of nonopen orbits. C.35 exhibits an exact obstruction
to extending that particular median prescription over an
adjacent open orbit. C.36 constructs compatible contractions
on closures with two-dimensional proper operator image and
proves continuity on every finite union of those closures.
C.37 supplies the uniform three-word estimate and gives a
continuous contraction on the closed union $X_2$ of all those
open orbits, including infinite accumulation. Extension across
$X_2$ to its exterior and frontier-compatible choices for the
dimension-zero and dimension-one operator images remain
additional requirements. C.12 retains its unresolved status.
